\documentclass[11pt]{article}
\usepackage{fontspec,amsmath,amssymb,amsthm,hyperref}
\usepackage[margin=1in]{geometry}
\setlength{\emergencystretch}{3em}
\title{The actual Gaussian norm probability obstruction}
\author{LangevinLean mathematical exposition}
\date{}
\begin{document}\maketitle
\noindent Source module: \texttt{AverageLMC/GaussianObstruction.lean}. Compile this proof document with XeLaTeX or LuaLaTeX. The accompanying HTML pairs the mathematical argument with the exact Lean statements and records the scope of the formalization.


This module uses the actual standard Gaussian measure in dimension p, with no assumed concentration event. It proves the exact variance of its squared norm and rules out a dimension independent order-p bound on its cubed norm with probability one half. This is the probability ingredient for the GLM counterexample.


\section{Compute the squared-norm variance from exact Gaussian moments}


The standard Gaussian second moment is \(p\), and the shared fourth-moment theorem gives \(p^2+2p\). Gaussian integrability of every finite moment proves square integrability of the squared norm before using the variance formula. Subtracting the squared mean gives \(\operatorname{Var}(\|Z\|^2)=2p\). Thus no finite-moment premise is supplied externally.



\noindent Lean declarations: \texttt{stdGaussian\_norm\_sq\_variance}.

\section{Apply Chebyshev to the small-norm event}


If \(\|Z\|^2<p/2\), its deviation from its mean p is at least p/2. Chebyshev bounds the probability by \(2p/(p/2)^2=8/p\). The proof uses a set inclusion and the actual Gaussian law, retaining the extended nonnegative probability values.



\noindent Lean declarations: \texttt{stdGaussian\_small\_norm}.

\section{Disprove every uniform cubic norm bound}


Let \(C\ge0\) and \(p>8C^2\). If \(\|Z\|^2\ge p/2\), then \(\|Z\|^6\ge p^3/8>C^2p^2\), so \(\|Z\|^3>Cp\). Therefore the event \(\|Z\|^3\le Cp\) is contained in the small-norm event. When also \(p>16\), its probability is at most \(8/p<1/2\). For every proposed C choose a natural dimension larger than both thresholds; its probability contradicts the proposed universal one-half lower bound.



\noindent Lean declarations: \texttt{stdGaussian\_cube\_bound\_fails}, \texttt{no\_uniform\_stdGaussian\_cube\_bound}.

\end{document}
